\documentclass{article}

\usepackage{amsmath, amssymb, amsthm}
\usepackage[a4paper, margin=1in]{geometry}

\title{Input Size vs Value}
\author{Sid}
\date{\today}

\newtheorem{theorem}{Theorem}
\newtheorem{lemma}{Lemma}

\begin{document}

\maketitle

% STATUS: done
% LAST EDITED:
% SLIDES DONE: no

\section{Goal}

Understand why trial division up to $\sqrt{n}$ is \textbf{not} polynomial time, even though it looks efficient at first glance.

\section{Core Idea}

Algorithms are analyzed in terms of \textbf{input size}, not the numeric value.

If the input is a number $n$, then input size is:
\[
    \text{input size} = \log n
\]

Trial division runs in:
\[
    O(\sqrt{n})
\]

We rewrite this in terms of input size:
\[
    \sqrt{n} = 2^{\frac{1}{2} \log n}
\]

\section{Intuition}

At first glance, $\sqrt{n}$ feels small compared to $n$.

But computers do not see $n$ directly. They see its binary representation.

So:
- $n$ grows exponentially with number of bits
- $\sqrt{n}$ is still exponential in number of bits

This means the algorithm is actually very slow.

\section{Example}

\subsection*{Example 1}

Given: $n = 10^{12}$

Steps:
- $\sqrt{n} = 10^6$
- So we do about 1 million checks

Now consider input size:
\[
    \log_2(10^{12}) \approx 40 \text{ bits}
\]

So:
- Input size = 40
- Work done = $10^6 \approx 2^{20}$

Result:
Work grows exponentially in input size.

\section{Proof}

\subsection*{Claim}

Trial division up to $\sqrt{n}$ is not polynomial in input size.

\subsection*{Proof Sketch}

Convert runtime from function of $n$ to function of $\log n$ and show it becomes exponential.

\subsection*{Proof}

\begin{proof}

    Let input size be:
    \[
        k = \log n
    \]

    Then:
    \[
        n = 2^k
    \]

    Trial division takes:
    \[
        \sqrt{n} = \sqrt{2^k}
    \]

    \begin{align*}
        \sqrt{n} & = (2^k)^{1/2} \\
                 & = 2^{k/2}
    \end{align*}

    Thus runtime:
    \[
        O(2^{k/2})
    \]

    This is exponential in $k$, not polynomial.

\end{proof}

\subsection*{Key Step}

The key step is rewriting $n = 2^k$ where $k$ is input size.
This exposes the exponential nature of $\sqrt{n}$.

\subsection*{Conclusion}

Trial division is exponential in input size, so it is not in P.

\section{Failure / Limitation}

The naive reasoning assumes:
- $\sqrt{n}$ is small compared to $n$

But ignores:
- input encoding size

This leads to incorrect conclusions about efficiency.

\section{Insight}

Efficiency must always be measured in terms of input size.

This reframes the problem:
- We need algorithms polynomial in $\log n$
- Not polynomial in $n$

\section{Algorithm (if applicable)}

\begin{enumerate}
    \item For $i = 2$ to $\sqrt{n}$
    \item Check if $i \mid n$
    \item If yes, composite
\end{enumerate}

\section{Complexity (if applicable)}

Time complexity:
\[
    O(\sqrt{n}) = O(2^{(\log n)/2})
\]

Space complexity:
\[
    O(1)
\]

\section{Connection}

This shows why naive methods fail and motivates faster algorithms like:
\begin{itemize}
    \item Fermat Test
    \item Miller-Rabin
    \item AKS
\end{itemize}

\section{Doubts}

Is there a deterministic algorithm that is polynomial in $\log n$?

\section{Sources}

\begin{itemize}
    \item Course notes
    \item Standard number theory references
\end{itemize}

\end{document}